\documentclass[11pt,a4paper]{article}

% ---------- Pacotes ----------
% Compila tanto com pdfLaTeX quanto com XeLaTeX/LuaLaTeX
\usepackage{iftex}
\ifPDFTeX
  \usepackage[T1]{fontenc}
  \usepackage[utf8]{inputenc}
  \usepackage[brazil]{babel}
  \IfFileExists{tgpagella.sty}{\usepackage{tgpagella}}{\usepackage{mathpazo}}
\else
  \usepackage{fontspec}
  \usepackage{polyglossia}
  \setdefaultlanguage[variant=brazilian]{portuguese}
  \setmainfont{TeX Gyre Pagella}
\fi
\usepackage[a4paper,margin=2.5cm]{geometry}
\usepackage{microtype}
\usepackage{setspace}
\usepackage{xcolor}
\usepackage{titlesec}
\usepackage{fancyhdr}
\usepackage{amsmath}
\usepackage{graphicx}
\usepackage{float}
\usepackage{caption}
\usepackage{xurl}
\usepackage[hidelinks]{hyperref}
\urlstyle{same}

\definecolor{destaque}{HTML}{1F4E79}
\setstretch{1.25}
\setlength{\parindent}{1.25cm}
\setlength{\parskip}{0.35em}

% Legendas das figuras
\captionsetup{font={small,stretch=1.1},labelfont={bf,color=destaque},labelsep=period}

% Títulos das questões
\titleformat{\section}[hang]{\normalfont\large\bfseries\color{destaque}}{}{0pt}{}
\titlespacing*{\section}{0pt}{1.6em}{0.8em}

% Cabeçalho e rodapé
\pagestyle{fancy}
\fancyhf{}
\fancyhead[L]{\small\color{gray}CFD Ambientes Oceânicos -- Atividade 01}
\fancyhead[R]{\small\color{gray}Breno Santos Cabral}
\fancyfoot[C]{\small\thepage}
\renewcommand{\headrulewidth}{0.4pt}

\begin{document}

% ---------- Identificação ----------
\thispagestyle{empty}
\begin{center}
  {\Large\bfseries\color{destaque} CFD Ambientes Oceânicos}\\[0.3em]
  {\large Atividade 03}\\[0.8em]
  {\normalsize Breno Santos Cabral}\\[0.2em]
  {\small 30 de outubro de 2026}
\end{center}
\noindent\textcolor{destaque}{\rule{\textwidth}{0.8pt}}

% ---------- Resultados ----------
O código \textsc{Fortran} fornecido foi reescrito na linguagem Python e pode ser acessado no link \url{https://github.com/BrenoSCabral/cfd_003}. As iterações considerando $\Delta t  = 1$ e $\Delta t = 10$ podem ser observada nas Figuras \ref{fig:dt1} e \ref{fig:dt10}, respectivamente. A diferença entre a solução analítica e as soluções numéricas pode ser observada no quadro inferior direito da Figura~\ref{fig:resultados}.

O código busca avaliar o desempenho de diferentes métodos numéricos na integração da equação

\begin{equation}
\frac{dU}{dt} = gU
\end{equation}

Para cada método numérico, a discretização está destacada a seguir:

\textbf{Progressivo}
\begin{equation}
U^{n+1} = g U^{n} 
\end{equation}

\begin{figure}[H]
  \centering
  \includegraphics[width=\textwidth]{../solucoes_dt1.png}
  \caption{Solução para $\Delta t = 1$ da equação de estudo de forma analítica e nos métodos Progressivo, Runge-Kutta Ordem 2 (RK2), Runge-Kutta Ordem 4 (RK4), Leap-frog, Matsuno, e Heun. O gráfico da esquerda apresenta os resultados em escala linear, enquanto o da direita utiliza escala logarítmica.}
  \label{fig:dt1}
\end{figure}

\begin{figure}[H]
  \centering
  \includegraphics[width=\textwidth]{../solucoes_dt10.png}
  \caption{Solução para $\Delta t = 10$ da equação de estudo de forma analítica e nos métodos Progressivo, RK2, RK4, Leap-frog, Matsuno, e Heun. O gráfico da esquerda apresenta os resultados em escala linear, enquanto o da direita utiliza escala logarítmica.}
  \label{fig:dt10}
\end{figure}



% ---------- Análise ----------
Utilizando o esquema de Euler progressivo, a equação $\dfrac{du}{dt} = F$ foi discretizada como
\begin{equation}
  u^{n+1} = u^{n} + F\,\Delta t .
\end{equation}

O método de von Neumann estabelece que a estabilidade pode ser avaliada a partir do fator de amplificação $\lambda$, definido por
\begin{equation}
  u^{n+1} = \lambda\, u^{n} ,
\end{equation}
sendo o esquema estável quando $|\lambda| \le 1$. Neste caso, como $F$ independe de $u$, $\lambda = 1$. Logo, o critério de estabilidade mostra que a solução é estável independentemente do passo de tempo adotado.

Além disso, como $dF/dt = 0$, o erro de truncamento é nulo. Por isso, mesmo com a diminuição da resolução temporal, a solução numérica permanece idêntica à solução analítica. De forma gráfica, uma iteração do esquema é representada na Figura~\ref{fig:iteracao}, para um $\Delta t$ qualquer, ilustrando a afirmação do parágrafo acima.



% ---------- Nota sobre IA ----------
\vspace{1.2em}
\noindent\fbox{\parbox{\dimexpr\textwidth-2\fboxsep-2\fboxrule}{\small\setstretch{1.1}%
\textbf{Uso de IA:} a ferramenta Claude (Anthropic) foi utilizada na revisão do texto, na elaboração da Figura~\ref{fig:iteracao} e na formatação em \LaTeX.}}

\end{document}
